\documentclass[aps,prd,preprint,superscriptaddress,amsmath,amssymb]{revtex4-2}
\usepackage{graphicx}   % Include figure files
\usepackage{bm}         % bold math
\usepackage{dcolumn}    % align table columns on decimal point
\usepackage{amsthm}
\usepackage{booktabs}
\usepackage{color}

\newtheorem{postulate}{Postulate}[section]
\newtheorem{theorem}{Theorem}[section]

\begin{document}

\title{The Geometric Description of Physics at the Cosmic Scale:\\ Mass, Information, Causality, Force, and the Quantum Relations from Geometry}

\author{Yuan-Jie Liu}
\email{yuanjieliu65@gmail.com}

\date{\today}

\begin{abstract}
We present a unified physical framework whose primitive is \emph{flowing
space}: space itself is a flowing medium, and all physical objects are
excitations (anchored patterns) of this flow.  From a small set of
postulates --- vector light speed (direction-changeable), three anchored
directions, and the counting of space-displacement ``strands'' --- the
framework derives, in a single chain of exact identities: mass as
symplectic entanglement volume, information as symplectic volume
conserved under unitary transport, and the four forces as the four
channels of the Leibniz decomposition of the flow momentum
$p=m(C-v)$.  The quantum relations $E=\hbar\omega$, the de~Broglie
relation $p=h/\lambda$, the Planck--Einstein relation $E=hf$, and the
uncertainty product $\Delta x\,\Delta p=\hbar$ emerge from helix
geometry, with the geometric choice $r=\lambda/2\pi$ (reduced wavelength
as helix radius); Pauli exclusion is the geometric rank criterion for
identical strands.  All identities are verified numerically with physical
constants at machine precision (relative errors $\lesssim 10^{-10}$), and
the derivation chain is fully formalized in Lean 4 with zero
\texttt{sorry}.  The honest status is
stated: the derivation \emph{chains are closed} (geometry $\Rightarrow$
quantum relations), while first-principle derivations of the geometric
choice $r=\lambda/2\pi$ and of the remaining input constants ($e$, $M_0$)
are not yet achieved.
\end{abstract}

\maketitle

\begin{quote}
\itshape
``Radicals always see matters in terms which are too simple. They divide
the universe into good and evil, black and white, us and them. But reality
is a spectrum, not a dichotomy.''
\par
\hfill --- Leto~II, \textit{God Emperor of Dune} (Frank Herbert, 1981)
\end{quote}

\section{Introduction}
\label{sec:intro}

Modern physics faces four conceptual contradictions that are not technical
difficulties but conflicts of axioms: the superposition of spacetime itself
(quantum superposition is linear; the general-relativistic background is
dynamical), the measurement problem (collapse is not in the Schr\"odinger
equation), the black-hole information paradox (unitarity versus the
horizon), and the vacuum-energy catastrophe ($10^{120}$ orders of
magnitude).  String theory and loop quantum gravity keep the quantum
axioms and add entities; they do not change the foundation.

The present work follows a different strategy: \emph{change the primitive}.
The primitive is not particles, not fields, not spacetime: it is
\emph{flow}.  Space is a flowing medium; particles and fields are
excitations --- anchored, spiralling patterns --- of this flow.  This
document collects the unified chain that has been built on this foundation \cite{ProjectionPhysics}
and formalized in the Lean 4 proof assistant (zero \texttt{sorry}), with numerical
verification (N1--N31) at machine precision.

\section{Postulates of the flowing-space framework}
\label{sec:postulates}

\begin{postulate}[Flowing space]
Space is a flowing medium.  A non-excited pattern is carried completely by
the flow (photon: $\mathrm d\tau^2=0$); an excited pattern is anchored,
deviating from the flow.
\end{postulate}

\begin{postulate}[Vector light speed]
The equivalent flow speed of space is the speed of light (the river model of black holes \cite{HamiltonLisle2008} is the same geometric structure); the vector light
speed $C$ has a direction that can change.  The flow momentum of a body is
$p = m(C-v)$, with $v$ the body velocity.
\end{postulate}

\begin{postulate}[Three directions]
Space is three-directional: $\sigma_1,\sigma_2,\sigma_3$ form an
irreducible, fully entangled basis ($(\sigma_1+\sigma_2+\sigma_3)^2=3I$).
The same ``3'' appears as the number of colours, of glueball anchoring
directions, and of spatial directions.
\end{postulate}

\begin{postulate}[Counting of strands]
Mass is the number of space-displacement strands moving at light speed
around the body: $m=\sqrt{N}\,M_0$ (anchored addition, $N$ strands).
Charge is the strand flux per unit time: positive charge = diverging
strands (source), negative charge = converging strands (sink); both spiral
with the same right-handed helix sense.
\end{postulate}

\section{Mass as symplectic entanglement volume}
\label{sec:mass}

With twistors $\pi_i$ as the half-spinor excitations \cite{Penrose1967}, the total
correlation matrix is $P = \sum_i \pi_i\pi_i^\dagger = AA^\dagger$ (the
columns of $A$ are the twistors).  The mass squared is the symplectic
entanglement volume:
\begin{equation}
m^2 = \det(AA^\dagger) = \sum_{S,\,|S|=n} \big|\det A_S\big|^2 ,
\label{eq:cauchy}
\end{equation}
the Cauchy--Binet chain: for $n=2$ and arbitrary $m$ (any number of
superposed half-spinors) the identity is fully formalized in Lean \cite{deMoura2015,mathlib2020}
($\det(\sum_i \pi_i\pi_i^\dagger)=\sum_{i<j}|\langle\pi_i,\pi_j\rangle|^2$,
GQS1); the general-$n$ expansion kernel (GQS2/3) and the fixed-subset
sum (GQS4) are likewise proved.  The unified chain is
\begin{center}
\begin{tabular}{lcc}
\toprule
Particle & $N$ & $m^2$ \\
\midrule
photon & 1 & $|\det_1|^2 = |\pi_1|^2 \cdot$ (rank-1, massless) \\
electron & 2 & $|\langle\pi_1,\pi_2\rangle|^2$ (symplectic area) \\
glueball & 3 & $|\det_3[\pi_1\pi_2\pi_3]|^2$ (volume) \\
\bottomrule
\end{tabular}
\end{center}
Excitation $=$ rank $=$ number of independent twistors: mass is
\emph{structural} (a rank criterion), not additive in the naive sense.
Figure~\ref{fig:triple} shows the numerical identity and the rank scan.

\begin{figure}[htbp]
\includegraphics[width=0.9\columnwidth]{figures_en/fig_triple_rank.png}
\caption{Left: $\det_3(P)$ versus $|\det_3[\pi_1\pi_2\pi_3]|^2$ (200 random
samples, exact identity).  Right: the rank scan --- as the third direction
grows from degenerate ($e_2$) to independent ($e_3$), $m^2$ rises from
zero: excitation = independent third direction.}
\label{fig:triple}
\end{figure}

\paragraph{Global coherence (QFT5).} The anti-phase identity
$\cos(\theta+\pi)=-\cos\theta$ holds everywhere: the correlation
$E(\theta)=\cos\theta\cos(\theta+\pi)$ is shape-invariant under the
propagation distance $L$ (all curves collapse onto $-\cos^2\theta$), with
$E(\Delta=\pi)=-1$ exactly.  Global entanglement is coherence inherent to
the space field itself --- no action at a distance; $E$ depends only on
the relative phase difference.  The right panel recalls the rank
criterion of Section~\ref{sec:mass}: for a two-twistor pair
$\det(p_1+p_2)=|\langle\pi_1,\pi_2\rangle|^2=\sin^2\alpha$, so the
excited mass is the relative orientation of the two half-spinors
(Fig.~\ref{fig:antiphase}).

\begin{figure}[htbp]
\includegraphics[width=0.9\columnwidth]{figures_en/fig_antiphase_global.png}
\caption{Global coherence and rank criterion: (left) the anti-phase
correlation $E=\cos\theta\cos(\theta+\pi)$ is identical for every
propagation distance $L$ ($E(\Delta=\pi)=-1$); (right) the two-twistor
determinant $\det(p_1+p_2)=\sin^2\alpha$: parallel spinors are massless,
orthogonal spinors give maximal mass.}
\label{fig:antiphase}
\end{figure}

\section{Information as symplectic volume}
\label{sec:info}

Information is redefined geometrically: \emph{information = symplectic
correlation volume} $\det(AA^\dagger)$ (rather than density-matrix
entropy).  Then:

\paragraph{Transport (GQC1).} The propagation operator of the flow $U$ is
unitary ($UU^\dagger=1$), and
\begin{equation}
\det\big((UA)(UA)^\dagger\big) = \det(AA^\dagger),
\end{equation}
with the transport formula $(UA)(UA)^\dagger = U(AA^\dagger)U^\dagger$.
Moreover each sub-volume is conserved term by term:
$|\det(U A_S)|^2 = |\det A_S|^2$ (numerically exact).  Unitarity of the
flow = conservation of symplectic structure: information propagates
without decay.

\paragraph{Causality (GQC2).} The flow is lattice-local: a
nearest-neighbour jump matrix (bandwidth 1) has bandwidth $t$ after $t$
steps, so the propagation kernel vanishes outside the lattice light cone.
Causality emerges from lattice locality (cf.\ Bell \cite{Bell1964} and Aspect \cite{Aspect1982} for the quantum side) --- the geometric version of the
no-communication theorem, without extra axioms.

\paragraph{Effective speed (GQC3).} In a non-uniform flow the equivalent
speed is $1+v_{\rm flow}$: superluminal \emph{equivalent} motion is a
geometric (coordinate-dragging) effect --- the
Alcubierre \cite{Alcubierre1994}/Painlev\'e--Gullstrand structure \cite{HamiltonLisle2008} --- while signals never exceed
the local speed of light.  The light cone is tilted by the flow
(Fig.~\ref{fig:cone}).

\begin{figure}[htbp]
\includegraphics[width=0.9\columnwidth]{figures_en/fig_tilted_lightcone.png}
\caption{Tilted light cone: with flow $v=1.5$ the causal region in the
rest frame has slope $1+v$ (equivalent speed $>1$), while in the flow
frame signals remain $\le 1$ per step.}
\label{fig:cone}
\end{figure}

\section{The four forces as the Leibniz decomposition of flow momentum}
\label{sec:forces}

With $p=m(C-v)$, the rate of change decomposes by the Leibniz rule into
four channels (GQF):
\begin{equation}
\frac{d}{dt}\big[m(C-v)\big]
= \dot m\,C + m\dot C - \dot m\,v - m\dot v ,
\label{eq:fourforces}
\end{equation}
identified with the electric ($\dot m C$; cf.\ Maxwell \cite{Maxwell1865}): mass flow $\times$ light
direction), nuclear ($m\dot C$: mass $\times$ light-speed flow --- change
of the vector light speed), magnetic ($\dot m v$: mass flow $\times$
velocity) and gravitational/inertial ($m\dot v$: mass $\times$
acceleration; cf.\ the geometric theory of gravity \cite{Einstein1915})
forces.  Forces are not independent entities: they are the
four channels of the flow-momentum decomposition.  Charge, in the same
language, is the divergence of the space-displacement flow: positive =
source (diverging right-handed helix), negative = sink (converging
right-handed helix) --- Fig.~\ref{fig:charge}.

\begin{figure}[htbp]
\includegraphics[width=0.9\columnwidth]{figures_en/fig_charge_helix.png}
\caption{Charge as strand flow: positive charge = diverging right-handed
helix (source); negative charge = converging right-handed helix (sink).
Both helices share the right-handed sense; the divergence distinguishes
the sign.}
\label{fig:charge}
\end{figure}

\section{The photon}
\label{sec:photon}

The photon is the excited state of an electron (cf.\ the light quantum \cite{Einstein1905}) in which mass and charge
have vanished: it is the helix structure left behind by the electron
entity.  Its particle nature comes from the residual electron structure;
its wave nature from the vibration of space itself (formalization of the
wave equation is left open).

Two models (Fig.~\ref{fig:photon}):
\begin{itemize}
\item \textbf{Model A}: one excited electron in a cylindrical helix; the
straight-line part moves at the speed of light.  This is a single twistor
--- rank 1, $\det=0$, massless (TW1).
\item \textbf{Model B}: two excited electrons rotating symmetrically
about a common axis, moving at light speed perpendicular to the rotation
plane.  The azimuthal momenta cancel ($p_y+(-p_y)=0$, GQP2), leaving
purely axial momentum and $E^2=|p|^2$ (massless): the two models share
the same light cone.
\item The observational correspondence of the two models (e.g.\ circular
polarizations) is left open.
\end{itemize}
The photon is at rest \emph{in} space and travels \emph{with} the flow
($v=C$: $p=m(C-v)=0$ with $m=0$, GQP3) --- consistent with the flow
postulate.

\begin{figure}[htbp]
\includegraphics[width=0.9\columnwidth]{figures_en/fig_photon_models.png}
\caption{Photon models.  Left: model A, single excited-electron helix
(axial speed $=c$).  Right: model B, two symmetric helices --- azimuthal
cancellation leaves purely axial momentum.}
\label{fig:photon}
\end{figure}

\section{The quantum relations emerge from helix geometry}
\label{sec:quantum}

\paragraph{Energy as rotation (GQR).}  The Planck constant $h$ is the
angular momentum per turn; the reduced Planck constant $\hbar=h/2\pi$ is
the per-radian angular momentum.  The frequency is the number of turns
per second (countable strands).  Then
\begin{equation}
E = \hbar\omega = \frac{h}{2\pi}\,2\pi f = h\,f =
\underbrace{h}_{\text{angular momentum per turn}}
\times
\underbrace{f}_{\text{turns per second}},
\end{equation}
the Planck--Einstein relation \cite{Planck1900} as the rotation energy of the helix, with the de~Broglie wavelength \cite{DeBroglie1924}
(Fig.~\ref{fig:turns}).

\begin{figure}[htbp]
\includegraphics[width=0.9\columnwidth]{figures_en/fig_photon_turns.png}
\caption{Counting the strands: each turn of the photon helix is one
frequency period and carries angular momentum $h$; $E=h\cdot f$ = angular
momentum per turn $\times$ turns per second.}
\label{fig:turns}
\end{figure}

\paragraph{Circumference = wavelength (GQR4-6).}  The factor $2\pi$
between $h$ and $\hbar$ is the circumference of the helix
cross-section.  With the geometric choice $r=\lambda/2\pi$ (the reduced
wavelength as helix radius), the chain closes:
\begin{align}
2\pi r &= \lambda && \text{(one turn = one wavelength),}\\
\hbar &= r\,p && \text{(angular momentum = radius $\times$ momentum),}\\
p &= \frac{h}{\lambda} && \text{(de~Broglie relation emerges),}\\
E &= pc = hf && \text{(Planck--Einstein emerges),}
\end{align}
verified with physical constants at machine precision
(Table~\ref{tab:circ}).

\begin{table}[htbp]
\caption{Circumference = wavelength: the derivation chain closed with
physical constants ($\lambda=500$\,nm).}
\label{tab:circ}
\begin{ruledtabular}
\begin{tabular}{lcc}
Quantity & Value & Check \\
\hline
$r=\lambda/2\pi$ & $79.58$\,nm & reduced wavelength \\
circumference $2\pi r$ & $500.0$\,nm & $=\lambda$ (one turn = one wavelength) \\
$p=h/\lambda$ & $1.3252\times10^{-27}$\,kg\,m\,s$^{-1}$ & de~Broglie \\
$\hbar=r\,p$ & $1.0546\times10^{-34}$\,J\,s & measured $1.05457\times10^{-34}$; rel.\ err.\ $6\times10^{-10}$ \\
$E=pc$ & $3.313\times10^{-19}$\,J & $=hf$ exactly \\
\end{tabular}
\end{ruledtabular}
\end{table}

\paragraph{Uncertainty from helix geometry (Heisenberg).}  The helix
radius is the natural scale of positional uncertainty,
$\Delta x \sim r = \lambda/2\pi$; the momentum is
$\Delta p \sim p = h/\lambda$.  Their product is
\begin{equation}
\Delta x\,\Delta p \;\sim\; \frac{\lambda}{2\pi}\,\frac{h}{\lambda}
= \frac{h}{2\pi} = \hbar \;\ge\; \frac{\hbar}{2},
\end{equation}
so the Heisenberg uncertainty relation
\cite{Heisenberg1927} is reproduced exactly at the natural scale of the
helix geometry (the helix cannot localize a photon better than its own radius
without destroying the strand structure).

\paragraph{Exclusion of identical strands (Pauli).}  Two identical
states are two parallel twistors; their symplectic volume vanishes,
$\det(AA^\dagger)=|\det A|^2=0$ --- identical strands create no new
structure, no mass, no new degree of freedom.  This is the geometric rank
criterion for exclusion: only \emph{distinct} strands (linearly
independent twistors) open a new channel, formally identical to the
geometric rank criterion behind the Pauli exclusion principle
\cite{Pauli1925}.

\section{Glueball, quark, and the one ``3''}
\label{sec:glueball}

The glueball mass rule reads as strand counting: $m=\sqrt{N}\,M_0$ with
$N=3$ anchored directions ($\sigma_1\sigma_2\sigma_3$); the lattice
$N$-sequence $\{3,6,7\}$ \cite{MorningstarPeardon1999,Chen2006,PDG2024} (channels $0^{++},2^{++},0^{-+}$) matches
$\sqrt{3},\sqrt{6},\sqrt{7}$ times $M_0$ within the lattice ranges.  The
number of colours is the same ``3'': colour $=3$ = three spatial
directions (structure correspondence: why 3 colours); 8 gluons
$= 2^3 = \dim C\ell(6)$ (TW5).  The ``3'' of space, of colour, and of
glueball anchoring is one and the same structure.

Honest limit: the quark mass spectrum ($2.2$\,MeV $\to$ $173$\,GeV)
shows no strand-counting pattern --- quark masses are QCD free
parameters; the colour count gives structure, not numbers.

\section{Numerical validation}
\label{sec:validation}

All identities above are verified numerically (N1--N31; Lean 4 \cite{deMoura2015,mathlib2020}, machine
precision): the Cauchy--Binet chain ($\le 5.9\times10^{-12}$),
unitary transport (formula error $0$, sub-volume conservation
$\le 5.9\times10^{-12}$), the light-cone bandwidth (violations $0$),
the four-force decomposition (error $0$), the photon model
(azimuthal cancellation $0$), the quantum matching
($E=\hbar\omega$ relative error $0$; $\hbar=rp$ relative error
$6\times10^{-10}$; $E=pc$ vs $E=hf$ error $0$), and the ``3'' chain
($\det_3=1$ for orthogonal unit directions, $m=\sqrt3\,M_0$).

\section{Honest assessment}
\label{sec:honest}

The following is claimed, and the following is not.

\paragraph{Claimed.} (i) A single self-consistent chain from one
primitive (flowing space) covers mass, information, causality, effective
speed, force, the photon, and the quantum relations; (ii) the derivation
chains are \emph{closed}: given the geometric choice $r=\lambda/2\pi$,
the de~Broglie and Planck--Einstein relations are consequences of helix
geometry plus classical identities; (iii) all statements are formalized in
Lean 4 with zero \texttt{sorry} and verified numerically at machine
precision.

\paragraph{Not claimed.} (i) A first-principle derivation of the
geometric choice $r=\lambda/2\pi$ (its origin is unknown); (ii) the
numerical derivation of the remaining constants: $G$ and $e$ are
\emph{structurally present} --- $G$ is the flow-gradient constant in the
rain speed $v(r)=\sqrt{2GM/r}$ (at the horizon $v=c$, exact), and $e$ is
the strand-flux quantum, with the fine-structure constant
$\alpha=e^2/(4\pi\varepsilon_0\hbar c)=1/137.036$ the geometric ratio of
$e$, $\hbar$ and $c$ (verified to $6\times10^{-10}$) --- but their values
still rest on the $M_0$ calibration and the measured $\alpha$; (iii) any
new falsifiable prediction beyond the standard framework (the lattice
$N$-sequence is a posteriori fitting); (iv) the continuum wave equation
for the photon's wave nature, and the divergence form of charge, remain
formalization skeletons.

\section{Conclusion}
\label{sec:conclusion}

The flowing-space framework unifies, in one chain of exact identities,
the concepts that physics has kept separate: mass (symplectic volume),
information (symplectic volume again), causality (lattice locality),
force (Leibniz channels of flow momentum), the photon (excited-electron
helix), and the quantum relations (rotation of the helix; circumference =
wavelength).  The chains are closed; the geometric choice at their head
is the open question.  Whether the foundation survives is a question for
experiment; the framework is now precise enough to be falsified.

\appendix
\section{The full derivation chain: from postulates to equations}
\label{app:derivations}

This appendix collects, in human-readable form, every derivation used in
the main text.  Each step is a complete algebraic identity; the
corresponding Lean~4 theorem \cite{deMoura2015,mathlib2020} and the
numerical verification (N1--N31) are cited step by step.  No step is
skipped; where a step requires an input that is \emph{not} derived, it is
declared explicitly.

\subsection{The postulates}
\label{app:postulates}

\begin{enumerate}
\item \textbf{Flowing space.} Space is a flowing medium.  A non-excited
pattern is carried completely by the flow (photon: $\mathrm d\tau^2=0$);
an excited pattern is anchored, deviating from the flow
(\texttt{unexcited\_comoving\_zero\_time}, \texttt{excited\_implies\_deviation\_from\_flow}).
\item \textbf{Vector light speed.} The equivalent flow velocity of space
is the speed of light, $|C|=c$, with $C$ direction-changeable.  The flow
momentum of a body is
\begin{equation}
p = m\,(C-v),
\label{eq:flowmomentum}
\end{equation}
where $v$ is the body's velocity (GQF1; \texttt{SLS1}).
\item \textbf{Three directions.} Space is three-directional:
$\sigma_1,\sigma_2,\sigma_3$ with
$(\sigma_1+\sigma_2+\sigma_3)^2=3I$ (three-direction entanglement basis,
QFT6; \texttt{three\_direction\_entanglement\_basis}).
\item \textbf{Strand counting.} Mass is the count of space-displacement
strands moving at light speed around the body, $m=\sqrt N\,M_0$
(anchored addition); charge is the strand flux per unit time (source =
positive, sink = negative; both right-handed helices).
\item \textbf{Geometric choice (input, not derived).} The helix radius
equals the reduced wavelength,
\begin{equation}
r = \frac{\lambda}{2\pi},
\label{eq:rchoice}
\end{equation}
declared as the one free geometric input of the quantum chain.
\end{enumerate}

\subsection{Chain 1: mass as symplectic entanglement volume}
\label{app:chain1}

With twistors $\pi_i$ as half-spinor excitations \cite{Penrose1967}, the
total correlation matrix is $P=\sum_i\pi_i\pi_i^\dagger=AA^\dagger$ (the
columns of $A$ are the twistors).  The mass squared is the symplectic
entanglement volume:
\begin{equation}
m^2=\det(AA^\dagger)=\sum_{S,\;|S|=n}\big|\det A_S\big|^2,
\label{eq:cb}
\end{equation}
the Cauchy--Binet identity.  Derivation steps (each formalized):
\begin{itemize}
\item $AA^\dagger$ is Hermitian; $\det(AA^\dagger)=\det(A)\det(A^\dagger)$
(Leibniz formula; \texttt{outer\_sum\_det}, GQN2).
\item $\det(A^\dagger)=\overline{\det(A)}$ (transpose-conjugate;
\texttt{det\_conj\_mul\_sum\_injective}, GQS2c), hence
$\det(AA^\dagger)=|\det A|^2$ for square $A$.
\item For $m>n$ (more twistors than dimensions), expanding
$\det(AA^\dagger)$ by multilinearity in the rows gives one term per
$n$-subset $S$ of the $m$ strands: the minors sum
$\sum_S|\det A_S|^2$ (\texttt{det\_conj\_mul\_sum\_expand} GQS2a,
\texttt{sum\_enum\_minor} GQS4).
\item $n=2$, arbitrary $m$ (any number of superposed half-spinors):
\begin{equation}
\det\Big(\sum_{i}\pi_i\pi_i^\dagger\Big)=\sum_{i<j}\big|\langle\pi_i,\pi_j\rangle\big|^2,
\label{eq:cb2}
\end{equation}
every unordered pair contributes its symplectic area
(\texttt{outer\_sum\_m2\_det}, GQS1).
\end{itemize}
The unified chain (GQN5, \texttt{excitation\_rank\_unified}):
$N=1$ photon (rank 1, $|\det_1|^2$, massless), $N=2$ electron
($|\langle\pi_1,\pi_2\rangle|^2$, symplectic area), $N=3$ glueball
($|\det_3[\pi_1\pi_2\pi_3]|^2$, volume).  The rank criterion:
excitation $\Leftrightarrow$ linearly independent twistors
(\texttt{triple\_excited\_implies\_independent} GQ4a;
\texttt{excited\_implies\_independent\_N} GQN4a); degenerate (parallel)
twistors give $\det=0$, i.e.\ massless
(\texttt{pair\_massless\_iff\_parallel} QFT8;
\texttt{triple\_degenerate\_massless} GQ5).  Numerical verification:
N19--N20 (Cauchy--Binet, $8\times10^{-13}$), N6 (rank criterion).

\subsection{Chain 2: information, causality, effective speed}
\label{app:chain2}

\textbf{Information} is redefined as the symplectic correlation volume
$\det(AA^\dagger)$.  Under unitary transport
$A\mapsto UA$ ($UU^\dagger=1$):
\begin{equation}
\det\big((UA)(UA)^\dagger\big)
=\det\big(U(AA^\dagger)U^\dagger\big)
=\det(AA^\dagger)\,\big|\det U\big|^2
=\det(AA^\dagger),
\label{eq:unit}
\end{equation}
using $\det(XY)=\det X\det Y$ twice and $\det(U^\dagger)=\overline{\det U}$
(\texttt{unitary\_transport\_conserves\_det}, GQC1; sub-volumes term by
term, N22).  Information propagates without decay.

\textbf{Causality} emerges from lattice locality.  Let a
nearest-neighbour jump matrix have bandwidth $1$.  Two facts:
(i) composing operators of bandwidth $a$ and $b$ gives bandwidth
$\le a+b$ (triangle inequality; \texttt{band\_comp}, GQC2b);
(ii) hence $t$ steps have bandwidth $\le t$ (\texttt{band\_pow}, GQC2b).
Therefore the propagation kernel vanishes outside the lattice light cone
(chain jump, \texttt{chainJump}, GQC2): no signal crosses the cone
(N23, violations $0$).

\textbf{Effective superluminality} is coordinate dragging.  In the flow
coordinates $\xi_k=k-\varphi_k$ a signal reachable in $t$ steps satisfies
$|j-i|\le t+|\varphi_j-\varphi_i|$ (\texttt{flowCoord}, GQC3a); with a
flow gradient $>1$ the rest-frame displacement exceeds $1$ in one step,
while the signal never exceeds the local light speed relative to the flow
(GQC3b).  The equivalent speed is $1+v_{\rm flow}$ exactly (N24).

\textbf{Global coherence} is field-inherent: the anti-phase identity
\begin{equation}
\cos(\theta+\pi)=-\cos\theta
\label{eq:antiphase}
\end{equation}
holds at every point, so the correlation
$E(\theta)=\cos\theta\cos(\theta+\pi)=-\cos^2\theta$ is independent of
the propagation distance $L$ and $E(\Delta=\pi)=-1$ exactly
(\texttt{antiphase\_anticorrelation\_global}, \texttt{antiphase\_correlation\_any\_point},
QFT5; N4).  No action at a distance: $E$ depends only on the relative
phase difference.

\subsection{Chain 3: the four forces}
\label{app:chain3}

From the flow momentum \eqref{eq:flowmomentum},
$p=m(C-v)$, the force is the Leibniz (product-rule) decomposition:
\begin{equation}
F=\frac{\mathrm dp}{\mathrm dt}
=\frac{\mathrm dm}{\mathrm dt}C
+m\frac{\mathrm dC}{\mathrm dt}
-\frac{\mathrm dm}{\mathrm dt}v
-m\frac{\mathrm dv}{\mathrm dt},
\label{eq:fourforces}
\end{equation}
four channels (GQF2):
\begin{center}
\begin{tabular}{ll}
\toprule
channel & physical identification \\
\midrule
$\dfrac{\mathrm dm}{\mathrm dt}\,C$ & electric force (strand flux; source = divergence) \\
$m\,\dfrac{\mathrm dC}{\mathrm dt}$ & nuclear force (direction change of the flow) \\
$-\dfrac{\mathrm dm}{\mathrm dt}\,v$ & magnetic force (flux against the body's velocity) \\
$-m\,\dfrac{\mathrm dv}{\mathrm dt}$ & gravitation (flow gradient) \\
\bottomrule
\end{tabular}
\end{center}
The identity is exact (product rule; \texttt{four\_force\_decomposition},
GQF2; N25, error $0$).  The identifications of the channels with the
four interactions are the interpretation layer; the divergence form of
charge and the wave equation are formalization skeletons.

\subsection{Chain 4: the photon}
\label{app:chain4}

The photon is the excited state of an electron \cite{Einstein1905} in
which mass and charge vanish.  Model B (two excited helices rotating
symmetrically about the axis): the azimuthal momenta cancel,
$p_y+(-p_y)=0$, so the net momentum is purely axial
(\texttt{azimuthal\_cancellation}, GQP2); the dispersion is lightlike,
$E^2=|p|^2$ (GQP2b).  Equivalently, at $v=C$ the flow momentum
\eqref{eq:flowmomentum} gives $p=m(C-v)=0$ (GQP3): the photon is at rest
in space, carried by the flow at light speed.  The two models share the
light cone (N26: azimuthal cancellation $0$, massless criterion
$\checkmark$).

\subsection{Chain 5: the quantum relations from helix geometry}
\label{app:chain5}

This is the central chain.  All steps are complete algebraic
identities; the single input is the geometric choice
\eqref{eq:rchoice}.

\begin{enumerate}
\item \textbf{Circumference = wavelength.}  With $r=\lambda/2\pi$,
the circumference of the helix cross-section is
\begin{equation}
2\pi r=2\pi\cdot\frac{\lambda}{2\pi}=\lambda,
\label{eq:circ}
\end{equation}
one turn of the helix unfolds exactly one wavelength
(\texttt{gqr4\_circumference\_wavelength}, GQR4).
\item \textbf{$\hbar$ as radius $\times$ momentum.}  The angular momentum
of one turn is
\begin{equation}
L=r\,p=\frac{\lambda}{2\pi}\,p,
\label{eq:hbar}
\end{equation}
the classical formula.  With $h=\lambda p$ (one quantum per turn,
$L=h$), \eqref{eq:hbar} reads
$\hbar=L/2\pi=h/2\pi$ --- Planck's constant is the angular momentum per
radian (\texttt{gqr5\_hbar\_eq\_rp}, GQR5; N29: $\hbar=rp$ matches the
true $\hbar$ to $6\times10^{-10}$).
\item \textbf{de~Broglie relation.}  From \eqref{eq:hbar},
\begin{equation}
p=\frac{h}{\lambda},
\label{eq:debroglie}
\end{equation}
the de~Broglie relation \cite{DeBroglie1924} emerges (GQR5).
\item \textbf{Planck--Einstein relation.}  With $E=\hbar\omega$ and
$\omega=2\pi f$:
\begin{equation}
E=\hbar\omega=\frac{h}{2\pi}\cdot 2\pi f=hf,
\label{eq:planck}
\end{equation}
and with \eqref{eq:debroglie}, $E=pc=(h/\lambda)c=hf$
(\texttt{gqr3\_planck\_relation}, \texttt{gqr6\_planck\_einstein},
GQR3/GQR6; N27: $E=\hbar\omega$ relative error $0$; $E=pc$ vs $hf$ error
$0$).  The quantum relations are consequences of helix geometry plus
classical identities, given \eqref{eq:rchoice}.
\item \textbf{Heisenberg uncertainty.}  The helix radius is the natural
scale of positional uncertainty, $\Delta x\sim r=\lambda/2\pi$, and
$\Delta p\sim p=h/\lambda$; hence
\begin{equation}
\Delta x\,\Delta p\sim\frac{\lambda}{2\pi}\cdot\frac{h}{\lambda}
=\frac{h}{2\pi}=\hbar\ge\frac{\hbar}{2},
\label{eq:heisenberg}
\end{equation}
the Heisenberg relation \cite{Heisenberg1927} at the natural scale of the
helix (the helix cannot localize a photon better than its own radius).
\item \textbf{Pauli exclusion.}  Two identical strands are two parallel
twistors; their symplectic volume vanishes,
$\det(AA^\dagger)=|\det A|^2=0$ --- identical strands open no new channel
(rank criterion of Chain~1).  Only distinct (linearly independent)
strands create a new degree of freedom, in the spirit of the Pauli
principle \cite{Pauli1925}.
\end{enumerate}

\subsection{Chain 6: gravitation and electrodynamics (structural)}
\label{app:chain6}

\textbf{Gravitation.}  A non-uniform flow $v(x)$ gives the Gordon metric
\begin{equation}
\mathrm d\tau^2=\mathrm dt^2-\frac{(\mathrm dx-v\,\mathrm dt)^2}{c^2},
\qquad g=\begin{pmatrix}1-v^2/c^2 & v/c^2\\ v/c^2 & -1/c^2\end{pmatrix},
\label{eq:gordon}
\end{equation}
with $g_{tt}=1-v^2/c^2$ matching the weak-field form
$g_{tt}=1-2\Phi/c^2$ at $\Phi=\tfrac12 v^2$: the gravitational potential
is half the square of the flow speed (\texttt{gordon\_weak\_field\_matches\_newton},
SG11).  The geodesic equation then gives the Newtonian limit
$a=-\nabla(\tfrac12 v^2)=-\nabla\Phi$ (verified numerically).  The
photon condition $|\mathrm dx-v\,\mathrm dt|=c\,\mathrm dt$ gives
$\mathrm d\tau^2=0$ (\texttt{gordon\_photon\_proper\_time\_zero}, SG9).
At the black-hole horizon $|v|=c$ and escape becomes impossible --- a
direct consequence of the postulate that the photon follows the flow
(PH2), not an extra axiom.

\textbf{Electrodynamics.}  With the space field $C$ as the only
independent field, define $E=-\partial_t C$, $B=\partial_x C$
(motion-kinematics).  Then Faraday's law is an identity,
$\partial_t B=\partial_t\partial_x C=\partial_x\partial_t C=-\partial_x E$
(\texttt{MS2}); Amp\`ere--Maxwell is equivalent to the wave equation of
$C$ (\texttt{MS3}).  Maxwell's equations reduce to the wave equation of
the space field; gauge redundancy is eliminated by the physical condition
$|C|=c$ (SLS1).  The 3D identities
$\nabla\cdot(\nabla\times C)=0$ and $\nabla\times(-\nabla\phi)=0$ are
automatic (SF1--SF2).

\subsection{Predictions}
\label{app:predictions}

Predictions of the framework, with honest status:

\begin{enumerate}
\item \textbf{Glueball lattice sequence} (a posteriori fit, verified):
$m=\sqrt N\,M_0$ with $N=\{3,6,7\}$ for $0^{++},2^{++},0^{-+}$
\cite{MorningstarPeardon1999,Chen2006,PDG2024}; $M_0\approx0.93\pm0.07$ GeV
consistent across the three channels.  Falsifiable by future lattice data
at $N=4,5$: $\sqrt4,\sqrt5\times M_0$.
\item \textbf{Three-direction timescale} (proposed; derivation and
formalization pending): the three-direction glueball
($m_G=\sqrt3\,M_0$) carries a timescale
\begin{equation}
\tau_G=\frac{1}{6M_0^2},
\label{eq:taug}
\end{equation}
i.e.\ $\approx0.17$\,GeV${}^{-1}\approx1.1\times10^{-24}$\,s at
$M_0\approx0.99$\,GeV.  The algebraic relation is recorded as a candidate
prediction; its physical identification (decay width, coherence time, or
other) and its derivation from the strand-counting postulate are open.
\item \textbf{Black-hole escape impossibility} (postulate consequence):
inside the horizon the photon follows the flow inward; outward modes
vanish (P2, P3 of the black-hole section).  Geometric, not an extra
axiom.
\item \textbf{No other new falsifiable predictions} are claimed beyond
the standard framework: the honest assessment of \S\ref{sec:honest}
stands.
\end{enumerate}

\subsection{Honest boundaries of the chain}
\label{app:boundaries}

\begin{center}
\begin{tabular}{lll}
\toprule
layer & examples & status \\
\midrule
mathematical identities & Cauchy--Binet, product rule, $\cos(\theta+\pi)$ &
true but standard \\
conceptual reframing & $\hbar$ = angular momentum per radian; force = four
channels; photon = helix & interpretation layer \\
derived inputs & \eqref{eq:rchoice}, $M_0$, $e$, $G$ values & inputs, not
derived \\
skeletons & wave equation; charge as divergence; continuum 3D &
formalization pending \\
\bottomrule
\end{tabular}
\end{center}

\bibliography{projection-physics}

\end{document}
